% Final proof: full independent root and cycles proof/source audits passed.
\paragraph{C.3as (Refuted): integral top finiteness in every positive characteristic and dimension.}
Let $k$ be an algebraically closed field of characteristic $p>0$ and
let $N\ge3$. There is an explicitly specified excellent normal strictly
henselian local $k$-algebra $R$, of dimension $N$ and residue field $k$,
whose puncture $X$ is regular and for which
\[
 B^{N-1}(X):=
 \mathrm{CH}^{N-1}(X)/\mathrm{CH}^{N-1}(X)_{\mathrm{alg}}
                    \longrightarrow (k,+)
 \tag{C.3as.1}
\]
is surjective. Algebraic equivalence here means the ordinary relation
defined by actual families on fixed products. In particular, the
integral top quotient is not finite, even when $k=\overline{\mathbb F}_p$.
The exhibited quotient has exponent $p$. This proves no nonvanishing
after tensoring with $\mathbb Q$, and makes no characteristic-zero claim.

\begin{proof}
Write $d=N-1\ge2$. We first give the projective data, then construct an
integral cycle invariant. The latter uses the Milnor Bloch formula and
the elementary $d\log$ map. It needs neither a comparison with higher
relative Quillen groups nor division by a factorial.

\emph{1. The complete tower data.}
Define integers
\[
 E_0=1,\qquad \ell_i=1+pE_{i-1},\qquad E_i=E_{i-1}\ell_i
               \quad(1\le i\le d-1),\qquad n=pE_{d-1}.
 \tag{C.3as.2}
\]
Every $\ell_i$ is greater than $p$ and prime to $p$.
In projective coordinates $(Y:Z:W)$ take
\[
 V_1:\quad Y^n-YW^{n-1}=Z^{n-1}W,\qquad
 \infty=(0:1:0),\qquad
 D_1=E_{d-1}(n-3)\infty,\qquad f_1=Z/W .
 \tag{C.3as.3}
\]
Here $n\ge6$. The affine equation $y^n-y=z^{n-1}$ has
$y$ derivative $-1$; at its unique point at infinity the $W$
derivative is $-1$. Thus $V_1$ is smooth and $z$ is an etale
coordinate on its affine chart. Since a smooth plane curve of
degree $n$ has canonical degree $n(n-3)$,
\[
                  \operatorname{div}(df_1)=n(n-3)\infty=pD_1.
 \tag{C.3as.4}
\]
The normalized differential has no residual zero. More explicitly,
with $q=Y/Z$ and $h=W/Z$,
\[
 h=q^n-qh^{n-1},\qquad
 z-q^{-n}=q^{\,n(n-3)}c,\qquad
 c=q(h/q^n)^{n-2}.
 \tag{C.3as.5}
\]
The function $q$ is a uniformizer, $h/q^n$ is a unit with value
one, and $c$ is a uniformizer. The term $q^{-n}$ is a $p$th
power, so (C.3as.5) also checks the exact normalized differential.

Inductively, suppose the smooth projective connected variety $V_j$
comes with an effective ample divisor $D_j$ and a rational function
$f_j$ such that $df_j$ is a nowhere vanishing section of
$\Omega^1_{V_j/k}(-pD_j)$. Put
\[
          \ell=\ell_{d-j},\qquad D'=D_j/\ell
                         \quad(1\le j\le d-1).
 \tag{C.3as.6}
\]
The indicated division is integral; its verification is included below.
The exact Frobenius sequence and $df_j$ give a nonzero extension
\[
 0\longrightarrow\mathcal O_{V_j}(-D_j)
   \longrightarrow E\longrightarrow\mathcal O_{V_j}
   \longrightarrow0
 \tag{C.3as.7}
\]
whose Frobenius pullback splits. One can define $E$ as the pullback
of that exact sequence along the section represented by $df_j$,
so the construction does not require choosing an unspecified
extension class. The negative ample bundles have no sections;
hence the connecting class is nonzero and the Frobenius splitting
is unique.

Use the bundle of lines $P=\mathbb P_{V_j}(E)$, with projection
$\bar g$ and infinity section $S_\infty$ corresponding to the
subbundle in (C.3as.7). Its normal bundle is
$\mathcal O_{V_j}(D_j)$. The Frobenius splitting gives a smooth
integral divisor $G$, disjoint from $S_\infty$, with affine equation
$s^p=f_j$. It is smooth because the normalized differential is
nowhere zero. Fibre degree and restriction to $S_\infty$ give
the precise divisor identity
\[
                   G\sim pS_\infty-p\bar g^*D_j.
 \tag{C.3as.8}
\]
Indeed $G-pS_\infty$ is a pullback, and its restriction is
$\mathcal O(-pD_j)$ because $G$ is disjoint from $S_\infty$.
This minus sign is essential. The source-sign discrepancy already
documented in C.3an is retained here; we use the derived formula
(C.3as.8), not the conflicting display in the author preprint.

Let $m=\ell-p$, so $0<m<\ell$ and $(m,\ell)=1$, and put
\[
 M=\mathcal O_P(-S_\infty)\otimes
                       \bar g^*\mathcal O_{V_j}(pD').
 \tag{C.3as.9}
\]
Then $M^{-\ell}=\mathcal O_P(G+mS_\infty)$.
The next variety is the normalized cyclic cover
\[
 V_{j+1}=
 \operatorname{Spec}_P
 \left(\bigoplus_{a=0}^{\ell-1}
        M^a\bigl(\lfloor am/\ell\rfloor S_\infty\bigr)\right),
 \qquad T^\ell=s^p-f_j,
 \tag{C.3as.10}
\]
with multiplication induced by that branch section and the displayed
scalar normalization. Equivalently, it is the normalization of $P$
in this specified function field. The simple valuation along $G$
makes the field extension have degree $\ell$, even when $\ell$
is composite. Denote the composite map by $g:V_{j+1}\to V_j$
and its reduced infinity section by $S$. Define
\[
                    D_{j+1}=(\ell-1)S+g^*D',
                    \qquad f_{j+1}=T .
 \tag{C.3as.11}
\]
These are the iterated data of
\citet[Definition~2.1, Lemma~2.2, Proposition~2.3, and
Section~2.4, author pp.~7--8 and 11--12]{D_Mukai2011},
with the divisor signs independently checked. The local and
ampleness verifications needed here follow next.

\emph{2. Smoothness, normal bundle, ampleness, and differential order.}
For a local equation $u$ of $D'$, write the local normalization
of the exact differential as
\[
       f_j=a^p+u^{p\ell}c,\qquad dc\ \hbox{nowhere zero},
       \qquad s=u^\ell s_i+a,\qquad T=u^pT_i.
 \tag{C.3as.12}
\]
These local forms come from the Frobenius sequence used in
(C.3as.7); equivalently, a local primitive of the normalized exact
differential gives $c$. They give
$T_i^\ell=s_i^p-c$.
On this finite chart $T_i,s_i$, together with complementary
coordinates to $c$ on the base, are smooth coordinates. This
also proves smoothness along the branch over $G$.

At infinity set $x=1/s_i$. Choose $b$ with
$1\le b<\ell$ and $pb\equiv-1\pmod\ell$, and set
$a_0=(1+pb)/\ell$. The integral local parameter
$w=x^{a_0}T_i^b$ satisfies
\[
                    w^\ell=x(1-cx^p)^b.
 \tag{C.3as.13}
\]
The derivative of the right side with respect to $x$ is a unit
at $x=0$. This is a smooth normal chart for the normalization:
it has the same function field, is finite near the branch, and
is regular there. Thus $S$ is $w=0$, maps isomorphically to
$V_j$, and $V_{j+1}$ is smooth everywhere. The cover is finite
over $P$ and therefore projective. Each fibre of $g$ is the
integral compactification of $T^\ell=(s-c_0)^p$; integrality
uses $(p,\ell)=1$.

On overlaps, the leading changes at infinity have the form
$x_j=(u_j/u_i)^\ell x_i+\text{higher terms}$ and
$T_j=(u_i/u_j)^pT_i$. Since $a_0\ell-pb=1$,
$w_j=(u_j/u_i)w_i$ to first order. Hence
\[
                         N_{S/V_{j+1}}=\mathcal O_{V_j}(D').
 \tag{C.3as.14}
\]
This also fixes any possible torsion ambiguity that would remain
if one used only the $\ell$th power of the normal bundle.

The effective divisor $S$ is nef: it has nonnegative intersection
with curves outside it, while on it the intersection is the degree
of the ample divisor $D'$. The divisor $A'=g^*D'$ is nef.
For an integral subvariety $W'\subset V_{j+1}$ of dimension $r>0$,
the fibres of $g$ have dimension one. If $\dim g(W')=r$,
then $(A')^r\cdot W'>0$. If $\dim g(W')=r-1$,
then $S(A')^{r-1}\cdot W'>0$, because the generic fibre is
integral and the section meets it with degree one. All the
other mixed intersections of these nef divisors are nonnegative.
Expanding $((\ell-1)S+A')^r\cdot W'$ proves it is positive
in both cases. Nakai--Moishezon therefore proves $D_{j+1}$
ample, with no assumption on the dimension.

From $T_i^\ell=x^{-p}-c$ we obtain
\[
                    \ell T_i^{\ell-1}dT_i=-dc .
 \tag{C.3as.15}
\]
At $S$, the order of $T_i$ is $-p$ and $dc$ is a nonzero
coordinate differential. Thus $dT_i$ has exact order
$p(\ell-1)$ along $S$. On the finite chart $dT_i$ is a
coordinate differential and has no residual zero. Since
$dT=u^p dT_i$, the entire divisor is
$p(\ell-1)S+pg^*D'=pD_{j+1}$, with no additional zero
at an intersection of these divisors. This proves the inductive
condition on the normalized differential.

Finally, $D_j$ is divisible by $E_{d-j}$. Initially this is
(C.3as.3). If it holds at stage $j$, then $D'=D_j/\ell_{d-j}$
is divisible by $E_{d-j-1}$ and
$\ell_{d-j}-1=pE_{d-j-1}$. Formula (C.3as.11) proves the
claim at stage $j+1$. Every step (C.3as.6) is consequently
defined. At the end put
\[
               V=V_d,\qquad L=\mathcal O_V(pD_d),
               \qquad
               \eta=df_d\in H^0(V,\Omega^1_{V/k}\otimes L^{-1}).
 \tag{C.3as.16}
\]
The variety $V$ is smooth connected projective of dimension $d$,
$L$ is ample, and $\eta$ is nowhere zero.

\emph{3. The specified strict-local cone.}
Take the full section ring and the strict henselization with
the indicated residue field:
\[
 A=\bigoplus_{\nu\ge0}H^0(V,L^\nu),\qquad
 R=(A_{A_{>0}})^{\mathrm{sh}},\qquad
 X=\operatorname{Spec}R\setminus\{\mathfrak m_R\}.
 \tag{C.3as.17}
\]
The ring $A$ is normal and finitely generated. The total space
$Z_0=\operatorname{Tot}(L^{-1})$ contracts its zero section by
a proper map to $\operatorname{Spec}A$, which is an isomorphism
away from the vertex. The argument from C.3an applies in any
dimension: choose a sufficiently large projectively normal
Veronese, use its usual cone contraction, and the finite fibre
power map. The resulting map to the full section ring is proper.
Sections in two consecutive sufficiently large powers recover
the degree-one fibre ratio locally, proving the asserted
isomorphism off the zero section. No generation of $A$ in
degree one is assumed.

Set $Z=Z_0\times_{\operatorname{Spec}A}\operatorname{Spec}R$.
It is regular, is proper over $R$, and has complement $X$ to
its exceptional divisor $i:V\hookrightarrow Z$. The projection
$\pi:Z\to V$ is flat: it is the composite of the flat
localization/ind-etale map $Z\to Z_0$ and the line-bundle
projection. The ring $R$ has dimension $d+1$ and all the
properties asserted. Its closed fibre in $Z$ has support $V$;
the divisor $V$ itself and its finite infinitesimal neighborhoods
are unchanged by this pointed etale base change. Indeed each such
neighborhood factors over an Artinian quotient of the local ring
at the vertex, whose pointed strict henselization has the same
quotient. This does not require the vertex ideal to pull back to
the ideal of the reduced divisor.

\emph{4. An integral conormal cycle invariant in all degrees.}
The scheme $Z$ is regular, noetherian, and contains the
infinite field $k$. Thus
\citet[Theorems~7.1 and 7.5, author pp.~21--22]{cyKerz2009Milnor}
applies integrally:
\[
             \mathrm{CH}^d(Z)=H^d(Z,\mathcal K_d^M).
 \tag{C.3as.18}
\]
The symbol map
$\{u_1,\ldots,u_d\}\mapsto d\log u_1\wedge\cdots\wedge d\log u_d$
is a morphism $\mathcal K_d^M\to\Omega^d_{Z/k}$ in every
characteristic. Products and residues have the usual integral
compatibilities; see
\citet[Sections~3--4, pp.~557--560, especially the residue
squares and equation~(24)]{cyBraunling2013}.

The retraction $\pi$ gives the canonical splitting
\[
 i^*\Omega^1_{Z/k}=\Omega^1_{V/k}\oplus L,\qquad
 i^*\Omega^d_{Z/k}=\Omega^d_{V/k}
                         \oplus(L\otimes\Omega^{d-1}_{V/k}).
 \tag{C.3as.19}
\]
Here $i^*$ tensors the sheaf with $\mathcal O_V$; it does
not apply the differentials map to $\Omega^1_{V/k}$ that
would kill the conormal summand. In a line-bundle chart that
summand is represented by $dt$ at $t=0$. Formula (C.3as.19)
is unchanged by the ind-etale map $Z\to Z_0$.
Apply (C.3as.18), $d\log$, restriction to
$i^*\Omega^d_{Z/k}$, and projection to this summand:
\[
 \delta:\mathrm{CH}^d(Z)\longrightarrow
                     H^d(V,L\otimes\Omega^{d-1}_{V/k}).
 \tag{C.3as.20}
\]
All these are maps on the actual schemes. In particular there
is no assertion that a formal class algebraizes. The local
computations may alternatively be performed on a pointed
smooth etale stage of $Z\to Z_0$; the Gersten descriptions,
symbols and coherent restrictions are compatible with their
filtered limit.

We check the entire exceptional image. On the actual total
space $Z_0$, the zero-section identity and projection formula
give, for $\beta\in\mathrm{CH}^{d-1}(V)$,
\[
                  i_{0*}\beta
                    =\pi_0^*(-c_1(L)\cdot\beta).
 \tag{C.3as.21}
\]
This follows by applying $i_{0*}i_0^*$ to $\pi_0^*\beta$
and using $\mathcal O_{Z_0}(V)=\pi_0^*L^{-1}$.
The square with the selected reduced zero section is
Cartesian after base change to $R$: on that section it is
the base change of $\operatorname{Spec}k$ to the pointed
strict henselization, and the result is still $V$.
Flat base change therefore yields the literal Chow identity
\[
                    i_*\beta=\pi^*(-c_1(L)\cdot\beta)
                                \quad\hbox{in }\mathrm{CH}^d(Z).
 \tag{C.3as.22}
\]
No homotopy invariance for $Z$ is being assumed. A class
pulled back from $V$ has zero conormal component in its
$d\log$ class, by (C.3as.19) and the product compatibility.
Consequently $\delta i_*=0$. The localization sequence
now descends (C.3as.20) to $\mathrm{CH}^d(X)$.

Pair with $\eta$ and use the Serre trace, fixing the order
of wedge factors once and for all:
\[
 \rho:\mathrm{CH}^d(X)\longrightarrow k,\qquad
 \rho(\xi)=\operatorname{Tr}_{V/k}\bigl(\delta(\xi)\wedge\eta\bigr).
 \tag{C.3as.23}
\]
The trace normalization is the functorial one in
\citet[Section~48.27]{D_Stacks0FVU}.

\emph{5. Ordinary algebraic equivalence is killed.}
The subgroup generated by actual ordinary curve-family
differences on $X$ is divisible, including by $p$.
Indeed extend a parameter curve to a smooth proper curve
$Q$ and close each horizontal cycle in $X\times Q$.
Horizontal integral closures are torsion free and hence
flat over the regular curve. They give a homomorphism
$\mathrm{CH}_0(Q)^0=\operatorname{Pic}^0(Q)(k)\to
\mathrm{CH}^d(X)$ by proper pushforward from $X\times Q$
after taking fibres. Principal divisors give the required
relations. The group of $k$-points of an abelian variety
over an algebraically closed field is divisible: each
multiplication map is a finite surjective map, also for
multiplication by $p$, and is surjective on $k$-points.
Finite sums of the resulting images are again divisible.
This proves the assertion for the ordinary subgroup.
Every homomorphism from that group to the exponent-$p$
group $(k,+)$ is zero. Thus (C.3as.23) factors through
the ordinary quotient in (C.3as.1).

\emph{6. Actual curves with all prescribed values.}
The construction contains a specified affine chart.
Starting with $y^n-y=f_1^{n-1}$, its successive equations are
\[
              f_j=s_j^p-f_{j+1}^{\,\ell_{d-j}}
                      \quad(1\le j\le d-1).
 \tag{C.3as.24}
\]
Eliminating $f_1,\ldots,f_{d-1}$ makes this chart finite
etale over $\mathbb A^d$ with coordinates
$s_1,\ldots,s_{d-1},f_d$: the remaining equation has
$y$ derivative $-1$. Let $v$ be the point at which all
these coordinates and $y$ are zero. Every $D_j$ is outside
the displayed chart in its new boundary directions, so
$L$ is trivial there and the normalized section $\eta$
is $df_d$. Use fibre coordinate $t$ on $L^{-1}$.

For each scalar $a\in k$ impose the actual polynomial
equations
\[
                      s_1=\cdots=s_{d-1}=0,\qquad f_d=at
 \tag{C.3as.25}
\]
and the remaining equation for $y$ obtained from (C.3as.24).
Its derivative in $y$ is $-1$. The point $y=t=0$ is
smooth and belongs to a unique irreducible component.
Take that component, and then its reduced closure in
$Z_0$. It is a curve not contained in the exceptional
divisor. Its proper map to its affine cone image is
birational, because the cone map is an isomorphism off
the exceptional divisor. Its vertex fibre is finite.
It is therefore the finite normalization of that image
after normalizing the curve.

More explicitly, normalization is finite by excellence,
and its proper map to the image is quasi-finite, hence
finite. After strict henselization at the image vertex,
this finite algebra splits into local factors. Select
the factor whose closed point is the specified germ $v$.
The curve was already smooth at that germ. Its normalization
is therefore an isomorphism there. Any remaining locus
where it is not an isomorphism is closed and proper over
this local one-dimensional branch; if nonempty, its closed
image would contain the closed point. This is impossible.
The chosen factor consequently embeds as an actual closed
integral curve $D_a\subset Z$ meeting $V$ only at $v$.
Its local ideal is exactly (C.3as.25).
This construction is done for each scalar individually;
no family on a varying henselized product is substituted
for a family on a fixed product.

\emph{7. The integral local residue, with coefficient one.}
For a regular sequence $g_1,\ldots,g_d$ cutting a curve
transversely at the point under consideration, its supported
Milnor cycle class is the successive boundary of
$\{g_1,\ldots,g_d\}$. Successive Gersten residues give
the multiplicity one of that complete intersection.
Applying $d\log$ gives the supported Hodge class
\[
                 \left[\frac{dg_1\wedge\cdots\wedge dg_d}
                                   {g_1\cdots g_d}\right].
 \tag{C.3as.26}
\]
This follows directly by multiplying the supported
principal-divisor classes and using the product/residue
compatibilities cited in Step~4. It is an integral identity,
not a Chern character divided by $(d-1)!$.

Apply it to
$(s_1,\ldots,s_{d-1},f_d-at)$ and subtract the class for
$a=0$. Restrict to $t=0$ while retaining $dt$ in
(C.3as.19). The ordinary $\Omega^d_V$ terms cancel and
the conormal part is
\[
 \pm a\,dt\otimes
       \left[\frac{ds_1\wedge\cdots\wedge ds_{d-1}}
                         {s_1\cdots s_{d-1}f_d}\right].
 \tag{C.3as.27}
\]
Pairing with the normalized $\eta=dt^{-1}\otimes df_d$
gives the local top residue $\pm a$. Its image under
the global Serre trace is the same number. One way to
fix this normalization is that the Koszul generator
in $\operatorname{Ext}^d(\mathcal O_v,\omega_V)$ maps
to $H^d(V,\omega_V)$ by the map dual to
$H^0(V,\mathcal O_V)\to H^0(v,\mathcal O_v)$, which
is the identity of $k$. Changing the fixed wedge
orientation if necessary, we obtain
\[
                        \rho([D_a]-[D_0])=a .
 \tag{C.3as.28}
\]
There is no contribution from another point, because
the actual selected curves meet $V$ only at $v$.
All scalar values are attained, proving (C.3as.1).
\end{proof}

The construction refutes integral top finiteness for every $p>0$
and every $N\ge3$. It leaves the unrestricted rational top quotient
and the unrestricted characteristic-zero top quotient undecided.
In particular no conclusion about a complex top quotient is drawn
from the positive-characteristic additive invariant.
